\chapter{Introduction}
\label{ch:introduction}

\section{Industrial Monitoring and the Motivation for Multimodal Learning}
Robotic welding motivated a practical question. If a camera records the appearance of the weld pool while other instruments record sound or process signals, when does a joint prediction improve on the strongest single-input model? Visual and acoustic welding studies have already implemented learned fusion \citep{P003,Hua2026Weld}. In one weld-penetration study, the image-only, audio-only, and fused models all achieved 100\% accuracy on the reported clean test set \citep[Tables~1--3]{Hua2026Weld}. For that comparison, the extra channel left no measurable margin for a clean accuracy gain. Whether it could help under camera occlusion or acoustic interference is a separate question that requires observations under those conditions.

The registered title reflects the original Intel Welding proposal. In the implementations, \textit{physics-informed} means that fusion weights can depend on observed physical process context, such as the assembly action being performed. The models impose no differential-equation or conservation-law constraints. The gear study was already under way when the Intel Welding audit found that acquisition sessions spanned partitions and proposed test recordings had been examined during development. The empirical evaluation consequently centred on five adjacent industrial and robotic monitoring tasks. Independent welding validation requires fresh welds. In the core assembly comparison, the visual classifier used the frozen ResNet18-vd CNN backbone inside RT-DETR-R18vd. A separate visual comparison evaluated the backbone together with the detector's spatial hybrid Transformer encoder.\ev{Researcher's account of the scope agreed with the supervisor}{\EvidenceRef{E003}{sections 2.3 and 3.1}}\ev{Sequence of the gear study and renewed welding audit}{\EvidenceRef{E004}{sections 11--14}}

The question across applications was what each input contributed to the prediction. In the gear study, this first required a measurable indicator of condition. Robotic assembly then supplied labelled actions with paired video and machine signals for a direct comparison of fusion methods. Milling, manipulation, and human--robot handover provided further tests on different operations. Each comparison is interpreted against its own outcome, available inputs, and strongest single-input reference.

\section{Problem Statement}
A monitoring model requires a target with a physical or operational interpretation and observations available at the intended decision time. A statistic extracted from a post-run photograph can define a continuous target for gear-condition regression. A label assigned to a completed assembly action defines an outcome-classification task. A camera frame from the end of that action may contain decisive evidence that was unavailable while the robot was executing it. These differences determine both what a model can learn and where its predictions could be used.

Estimating the value of an additional input also requires matching the comparison. Visual and sensor observations must refer to the same event, and both the fused model and its component predictors must be assessed on the same held-out units. Changing the training budget, fitting a new calibration function, or exposing only one model to simulated failures changes the interpretation of a score difference. Known absence raises a different problem from an apparently valid sensor reading corrupted by noise. Prior work on multimodal learning has shown that a joint model can perform worse than its strongest component \citep{Wang2020What}. The useful comparison therefore includes strong single-input models, simple combinations, and explicit tests of input degradation.

\section{Research Aim and Objectives}
The aim is to determine when combining visual and machine-sensor evidence improves industrial condition estimation or outcome classification. The implementation provides interchangeable visual and sensor representations, unimodal prediction heads, and fusion at the feature or prediction level. The comparison includes convolutional image features, statistical summaries of sensor windows, and patch-based Transformer encoders for sensor sequences. Static fusion uses fixed rules or fitted coefficients. A process-aware gate can change the relative weights with the context of an observation.

The evaluation asks whether those choices improve ranking or regression error, whether predicted probabilities and decisions remain useful, and how each method behaves when an input is unavailable or degraded. Grouped assessment accounts for dependence between observations from a shared experiment, recording, or task where the necessary identifiers are available. The interpretation also considers costs of computation, uncertainty in the target, and the extent to which a result carries over to new operating conditions.

\section{Main Research Question and Four Subquestions}
The original expos\'{e} posed the following main research question.
\begin{quote}
How can a process-aware multimodal AI architecture be designed and evaluated for robust industrial anomaly detection using robotic welding data as a case study?
\end{quote}
Four subquestions address representation, fusion, robustness, and evaluation.
\begin{enumerate}[label=SQ\arabic*.,leftmargin=16mm]
\item Which lightweight or pretrained Transformer-based representation components are suitable for visual, audio, and sensor time-series data in a multimodal industrial anomaly detection pipeline?
\item How do unimodal baselines, simple concatenation fusion, modular late fusion, and process-aware gated cross-modal fusion compare in terms of predictive performance?
\item Can a gated cross-modal fusion architecture improve robustness when one or more modalities are missing, degraded, or unreliable?
\item How can multimodal architecture performance be evaluated beyond accuracy, including robustness, interpretability, modularity, and practical usability?
\end{enumerate}
The welding case motivated all four questions.\ev{Research questions and scope}{\EvidenceRef{E002}{pp. 2--3}, \EvidenceRef{E001}{}} Audio remained relevant to the literature review. For REASSEMBLE, the correspondence between action annotations and waveforms was unresolved, and spoken outcomes posed an additional risk of revealing the label. The predictive experiments therefore concentrated on vision and machine sensing.\ev{REASSEMBLE audio feasibility}{\EvidenceRef{E004}{section 15}}

\section{Scope and Evolution of the Research}
The PHM study asked what the gear measurements could establish about condition. Post-run photographs supplied a provisional continuous measure of candidate damage area. The main sensor comparison held out whole experiments and tested whether three groups of measurements could predict that measure more accurately than the training mean. Their mean absolute error was 0.824 percentage points, compared with 0.804 for the constant predictor (Table~\ref{tab:phm-loeo}). The paired analysis left the difference unresolved. An independent physical measurement would be needed to determine how faithfully the photographic target tracks damage.

REASSEMBLE supplied the principal fusion test because video and sensor streams could be linked to native success or failure labels for the same completed actions. Across 4,530 actions from 148 recordings, failure average precision was .429 for the statistical sensor predictor, .446 for equal averaging of visual and sensor probabilities, and .447 for logistic stacking (Tables~\ref{tab:reassemble-unimodal} and~\ref{tab:reassemble-fusion}). The paired comparison supported a modest gain for stacking over its sensor component. Its advantage over equal averaging remained uncertain. The gates conditioned on quality alone or on quality and action context showed no consistent additional benefit when streams were absent or degraded. The measured gain applies to the fitted pipeline, which also introduced a second probability model.

The external implementations examined whether the same approach remained useful on new targets. In MSM, camera and machine measurements covered milling intervals labelled for abnormal cutting. For the 1,895 intervals in the primary evaluation, average precision for abnormal cutting, averaged over three seeds, was .757 for stacking and .829 for logistic sensing (Table~\ref{tab:msm-clean}). In ViFailback, the target was the outcome of a completed manipulation episode, and entire tasks were held out. Original stacking had lower pooled macro average precision than the logistic sensor model at each of three seeds. A separate experiment transferred a convolutional encoder learned on REASSEMBLE and found mixed changes across seeds (Tables~\ref{tab:vif-bothclass} and~\ref{tab:vif-matched}).

HANDOVER examined whether video and force--torque signals from the same part of a human--robot handover improved prediction of its recorded outcome. On 151 official test trials, equal averaging achieved failure average precision of .803, compared with .824 for the logistic sensor predictor (Table~\ref{tab:handover-clean}). The participant-to-trial mapping needed to verify participant-disjoint predictions for fitting a learned fusion model was unavailable. The empirical result therefore concerns simple averaging on those trials. Together, the three external studies show why a new application of a fusion pipeline needs a fresh comparison with its own sensor baseline.

\section{Contributions of the Thesis}
The implementation specifies how the models transform each physical observation into a prediction. In REASSEMBLE, the pretrained RT-DETR-R18vd checkpoint supplied its own ResNet18-vd convolutional backbone. The implementation pooled its features and trained a binary action classifier. The separate PHM image regressor collected intermediate features during an RT-DETR-L prediction pass. The evaluation linked each fitted model to its training and assessment groups, calibration procedure, and component predictors. This makes explicit which comparisons test combining two predictions and which also test whether their relative weights should vary from one observation to the next. It also exposed limits, including a preprocessing-isolation concern in a separate PHM four-arm comparison.

The second contribution is an account of reliability that distinguishes a known missing input, a degraded but present input, and a change of operating conditions. A programmed fallback can select the surviving sensor stream once the missing input has been identified. Recognizing an unreliable reading that still looks valid requires an additional form of evidence. Controlled tests evaluated input loss and corruption. The available group comparisons offered more limited evidence about changes in physical operating conditions. Later pooling, fine-tuning, and transfer studies are interpreted under their respective fitting procedures and degree of prior assessment exposure.

The literature review examined \ReviewWorkCount{} works in full text across two stages of screening and appraisal. Its comparison matrix records the task, physical input sources, fusion method, unimodal controls, assessment design, tests of unreliable inputs, and source locations. The welding comparison illustrates why these controls matter. A fused model can achieve high accuracy even when a single input yields the same result. The empirical chapters examine how much the combined model adds on each task and, where degradation tests are available, whether that gain persists.

\section{Organization of the Thesis}
Chapter~2 compares the closest welding and industrial fusion studies. Chapters~3 and~4 define the model components and evaluation procedures. Chapter~5 examines construction and prediction of the gear-condition proxy, and Chapter~6 presents the main assembly comparison. Chapters~7 and~8 report the milling and manipulation studies. Appendix~\ref{app:handover} gives the supplementary handover evaluation. Chapter~\ref{ch:discussion} draws on all five studies to answer the four subquestions and assess how far their findings answer the original welding question. Chapter~\ref{ch:conclusion} summarizes the contribution and the evidence still required for weld-quality assessment. The remaining appendices document the literature appraisal, datasets, implementation, sensitivity analyses, human review, reproducibility, research chronology, and prompt records.
